\section{Surface invariants and the cyclic cover}
\label{sec:surfaces}

Throughout, a curve means a reduced irreducible compact analytic curve unless
another meaning is specified. A rational curve is a curve whose normalization
is $\PP^1$; the curve itself is allowed to be singular. We orient every complex
surface by its complex structure, and write a dot for the intersection pairing.
We use $K_X$ both for the canonical line bundle and for its first Chern class
when a cohomological operation makes the meaning unambiguous.

The compact-surface identities of Kodaira give, under the hypotheses of
Theorem~\ref{thm:main},
\begin{equation}\label{surface:invariants}
 h^{0,1}(X)=1,\qquad h^{1,0}(X)=h^{2,0}(X)=0,\qquad
 b_2^+(X)=0,\qquad \chi(\OO_X)=0,
 \qquad K_X^2=-b_2(X).
\end{equation}
These are the numerical identities for class VII surfaces
\cite[Theorem~3]{Kodaira1964}; see also \cite{BHPV2004}.
In particular the intersection form on $H^2(X;\RR)$ is negative definite.
Our two surface-theoretic reductions are the following established results.

\begin{proposition}[Surface reductions]\label{surface:reduction}
Let $V$ be a minimal compact complex surface with $b_1(V)=1$,
$\kappa(V)=-\infty$, and $b_2(V)>0$.
\begin{enumerate}[label=\textup{(\roman*)}]
\item If $V$ has a nonzero effective divisor $C$ with $C^2=0$, then $V$
has a global spherical shell.
\item If $V$ has exactly $b_2(V)$ rational curves, then $V$ has a global
spherical shell.
\end{enumerate}
Here the conclusion supplies an embedded open neighborhood of
$S^3\subset\CC^2\setminus\{0\}$ with connected complement in $V$.
\end{proposition}

The first assertion is the square-zero-divisor case of Enoki's theorem
\cite{Enoki1981}, stated in this form in \cite[p.~1758]{Teleman2010}.
The second is the Main Theorem of Dloussky--Oeljeklaus--Toma
\cite[p.~283]{DOT2003}; their definition on p.~284 includes the
connected-complement condition. These statements impose no condition on the
fundamental group beyond the displayed surface hypotheses. The original
surface $X$ satisfies all of them. We may therefore assume henceforth that
\begin{equation}\label{surface:no-square-zero}
 C^2<0\quad\text{for every nonzero effective divisor }C\text{ on }X.
\end{equation}
Indeed negative definiteness leaves only the already settled square-zero
case as an alternative.

\begin{lemma}[Canonical basis and curve degrees]\label{surface:degree}
Put $b=b_2(X)$. There is a basis $u_1,\ldots,u_b$ of
$H^2(X;\ZZ)/\operatorname{Tors}$ with
\begin{equation}\label{surface:canonical-basis}
 u_i\cdot u_j=-\delta_{ij},\qquad
 K_X=u_1+\cdots+u_b\pmod{\operatorname{Tors}}.
\end{equation}
For every curve $B\subset X$,
\begin{equation}\label{surface:curve-degrees}
 B^2<0,\qquad (K_X+B)\cdot B\leq0,\qquad K_X\cdot B\geq0.
\end{equation}
Consequently every effective divisor has nonnegative canonical degree.
\end{lemma}

\begin{proof}
Donaldson's diagonalization theorem in its arbitrary-fundamental-group
version applies to the smooth, closed, oriented four-manifold $X$, whose
intersection form is negative definite by \eqref{surface:invariants}.
Precisely, we use Theorem~1 of \cite{Donaldson1987}, on integral cohomology
modulo torsion. Its surface application is also recorded in
\cite[\S1.1, p.~1756]{Teleman2010}.
Choose a negative orthonormal integral basis and write
$K_X=\sum_i k_i u_i$ modulo torsion. The canonical class is characteristic:
its mod-two reduction is the second Stiefel--Whitney class. Thus every $k_i$
is odd. Since $K_X^2=-b$, we have $\sum_i k_i^2=b$, so $k_i=\pm1$ for
every $i$. Replacing each basis vector by its negative when necessary gives
\eqref{surface:canonical-basis}.

If $[B]=\sum_i a_i u_i$ modulo torsion, the integers $a_i$ give
\[
 (K_X+B)\cdot B=-\sum_i a_i(a_i+1)\leq0.
\]
The strict inequality $B^2<0$ follows from
\eqref{surface:no-square-zero}. The adjunction formula for a possibly
singular integral curve is
\[
 K_X\cdot B=2p_a(B)-2-B^2.
\]
If $p_a(B)\geq1$, this is positive. If $p_a(B)=0$, the normalization-genus
formula $p_a(B)=g(\widetilde B)+\sum_x\delta_x$ shows that $B$ is a smooth
rational curve. Minimality and $B^2<0$ then give $B^2\leq-2$, so its
canonical degree is again nonnegative. Additivity proves the last assertion.
The adjunction and characteristic-class formulas used here are the standard
compact-surface formulas, as in \cite{BHPV2004}.
\end{proof}

Since $H^1(X;\ZZ)\simeq\ZZ$, choose a primitive generator $\alpha$.
The induced epimorphism $\pi_1(X)\to\ZZ$ defines a connected regular
covering
\[
 \pi:Y\longrightarrow X,
 \qquad \operatorname{Deck}(Y/X)=\langle T\rangle\simeq\ZZ.
\]
Only this quotient of the fundamental group is used. Choose a smooth real
closed representative of $\alpha$ and integrate its pullback to obtain a
smooth function $h_0:Y\to\RR$ satisfying
$h_0(Ty)=h_0(y)+1$. Later we may reverse both the generator and the sign of
the height. The Hermitian metric on the underlying surface $Y$ will always
be lifted from $X$; auxiliary bundle metrics will be specified separately.

For $\lambda\in\CC^*$, let $P_\lambda$ be the flat holomorphic line bundle
whose sections pull back to holomorphic functions with
$f(Ty)=\lambda f(y)$. This convention fixes the sign of the character.
These bundles are smoothly trivial: if $e^u=\lambda$, the function
$e^{u h_0}$ is a nowhere-zero smooth equivariant section. In particular
$c_1(P_\lambda)=0$ integrally.

\begin{lemma}[Character cohomology]\label{surface:characters}
For $\lambda\neq1$ and every $q=0,1,2$,
\begin{equation}\label{surface:acyclic}
 H^q(X,P_\lambda)=0,
 \qquad H^q(X,K_X\otimes P_\lambda)=0.
\end{equation}
For every $\lambda\in\CC^*$, including $\lambda=1$,
\begin{equation}\label{surface:canonical-vanishing}
 H^0(X,K_X\otimes P_\lambda)=0.
\end{equation}
At the trivial character,
$H^0(X,\OO_X)=\CC$, $H^1(X,\OO_X)=\CC$, and $H^2(X,\OO_X)=0$.
If a holomorphic function $f$ on $Y$ satisfies $f\circ T-f=c$ for a
constant $c$, then $f$ is constant and $c=0$. More generally, a holomorphic
function annihilated by a positive power of $T^*-1$ is constant.
\end{lemma}

\begin{proof}
A nonzero section of $P_\lambda$ has an effective zero divisor of Chern
class zero. A nonempty such divisor contradicts
\eqref{surface:no-square-zero}; hence the section is nowhere zero.
Its lifted function $f$ has a logarithmic differential $df/f$, defined
without a choice of logarithm. Equivariance makes $df/f$ invariant under
$T$, so it descends to a holomorphic one-form on $X$. By
\eqref{surface:invariants} this form vanishes. Thus $f$ is a nonzero
constant and $\lambda=1$. This proves the degree-zero assertion for
$P_\lambda$.

A nonzero section of $K_X\otimes P_\lambda$ would have zero divisor $C$
with $[C]=K_X$ in real cohomology. If $C$ is nonempty, then
$K_X\cdot C=K_X^2=-b<0$, contradicting Lemma~\ref{surface:degree}.
If $C$ is empty, the bundle is holomorphically trivial and its real first
Chern class vanishes, which also contradicts $K_X^2<0$. This proves
\eqref{surface:canonical-vanishing} for all characters.

Serre duality gives
$H^2(X,P_\lambda)\simeq
H^0(X,K_X\otimes P_{\lambda^{-1}})^*=0$.
Riemann--Roch for line bundles gives
$\chi(P_\lambda)=\chi(\OO_X)=0$, since $c_1(P_\lambda)=0$.
Thus $H^1(X,P_\lambda)=0$ if $\lambda\neq1$. Serre duality once more
gives all the canonical-bundle vanishings in \eqref{surface:acyclic}.
The assertions at $\lambda=1$ follow from
\eqref{surface:invariants} and connectedness of $X$.

Finally, $f\circ T-f=c$ makes $df$ invariant, hence the pullback of a
holomorphic one-form on $X$. It is zero, so $f$ is constant. For the more
general assertion, write $\Delta=T^*-1$. If a least exponent $m>1$ had
$\Delta^m f=0$ and $\Delta^{m-1}f\neq0$, the last function would descend
to a nonzero constant on $X$. The function $\Delta^{m-2}f$ would then
have a nonzero constant difference, which was just excluded. Thus $m=1$,
and $f$ descends to a constant on $X$.
\end{proof}

\begin{lemma}[A height with positive complex Laplacian]\label{surface:height}
There are a Gauduchon metric on $X$, a choice of $T$, and a smooth proper
height $h:Y\to\RR$ such that
\begin{equation}\label{surface:positive-height}
 h(Ty)=h(y)+1,\qquad Lh=c>0.
\end{equation}
Here $L$ is the complex Laplacian with the subharmonic sign. All positive
order derivatives of $h$ are uniformly bounded in the lifted metric.
\end{lemma}

\begin{proof}
Gauduchon's theorem supplies a Hermitian fundamental form $\omega$ with
$\partial\dbar\omega=0$ \cite{Gauduchon1977}. Fix the convention
$d^c=i(\dbar-\partial)$, so $dd^c=2i\partial\dbar$, and put
$L=\Lambda_\omega dd^c$. This is elliptic and has the subharmonic sign.
Its kernel on real functions on the compact connected surface consists
of constants by the maximum principle. Its index is zero: its principal
symbol is that of a positive multiple of the scalar Laplacian, and a
homotopy of its lower-order terms to that Laplacian preserves the elliptic
index. Moreover,
\[
 \int_X Lu\,\frac{\omega^2}{2}
 =\int_X dd^c u\wedge\omega
 =\int_X u\,dd^c\omega=0.
\]
The cokernel therefore consists exactly of constants. These are also the
Gauduchon-Laplacian solvability facts stated in
\cite[p.~289]{Buchdahl1999}.

The one-forms $dh_0$ and $d^ch_0$, and the function $Lh_0$, are periodic.
Let $c$ be the mean of $Lh_0$ on $X$. The preceding solvability yields a
smooth function $v$ on $X$ such that
$L(v\circ\pi)=c-Lh_0$. Set $h=h_0+v\circ\pi$.
If $c=0$, the descended real $(1,1)$-form $\eta=dd^ch$ would be both
exact and primitive. On a Hermitian surface a real primitive $(1,1)$-form
is anti-self-dual, and consequently
\[
 0=\int_X\eta\wedge\eta
   =-\int_X|\eta|^2\,\frac{\omega^2}{2}.
\]
The first equality follows from exactness and Stokes' theorem. Thus
$dd^ch=0$, so $\partial h$ is an invariant holomorphic one-form.
Its descended form is nonzero: otherwise the real function $h$ would be
constant, contrary to its unit increment under $T$. This contradicts
$h^{1,0}(X)=0$. Hence $c\neq0$; replacing $(T,h)$ by $(T^{-1},-h)$ if
necessary makes $c>0$.

Every positive order derivative of $h$ is periodic and therefore bounded.
To see properness, choose a compact set $F\subset Y$ whose deck translates
cover $Y$, possible because $X$ is compact. The height is bounded on $F$.
For any bounded interval $I$, only finitely many translates $T^jF$ can
meet $h^{-1}(I)$, by the unit-increment identity. The closed set
$h^{-1}(I)$ is therefore compact. This proves properness.
\end{proof}

\begin{lemma}[Finite cyclic quotients]\label{surface:covers}
For every $n\geq1$, put $S=T^n$ and $X_n=Y/\langle S\rangle$, and let
$p_n:X_n\to X$ be the degree-$n$ covering. Then $X_n$ is a minimal class
VII surface and
\[
 \chi(\OO_{X_n})=0,\qquad b_1(X_n)=1,\qquad
 b_2(X_n)=n b_2(X),\qquad K_{X_n}=p_n^*K_X.
\]
In particular its integral intersection lattice modulo torsion has a
negative orthonormal basis whose sum is $K_{X_n}$ modulo torsion.
\end{lemma}

\begin{proof}
The unramified holomorphic covering is a local biholomorphism, so its
tangent and canonical bundles are pullbacks. Compact-complex
Riemann--Roch, or the degree formula for the Todd class
\cite[\S4]{AtiyahSingerIII1968}, gives
$\chi(\OO_{X_n})=n\chi(\OO_X)=0$.
If $\sigma$ were a nonzero section of $K_{X_n}^{\otimes m}$ for $m>0$,
the product of its $n$ deck transforms would be a nonzero invariant
section of $K_{X_n}^{\otimes mn}$. It would descend to a nonzero section
of $K_X^{\otimes mn}$, contradicting $\kappa(X)=-\infty$.
Thus $\kappa(X_n)=-\infty$ and $p_g(X_n)=0$; the Euler characteristic
then gives $q(X_n)=1$.

The surface $X_n$ is non-K\"ahler: averaging a hypothetical K\"ahler form
over its finite deck group would give an invariant K\"ahler form and
hence a K\"ahler form on $X$. That is impossible because $b_1(X)=1$.
The compact-surface degree-one Hodge identities give
$b_1=2q$ for a surface with even $b_1$ and $b_1=2q-1$ for a surface with
odd $b_1$ \cite{BHPV2004}. By the even-first-Betti K\"ahler criterion
\cite[Theorem~11]{Buchdahl1999}, $b_1(X_n)$ is odd. Hence
$b_1(X_n)=2q(X_n)-1=1$.
The topological Euler characteristic of a closed oriented four-manifold
is $2-2b_1+b_2$. It multiplies by the degree of a finite covering, so
$b_2(X_n)=n b_2(X)$.

For minimality, suppose that $C\subset X_n$ were a smooth rational
curve with $C^2=-1$. Adjunction would give $K_{X_n}\cdot C=-1$.
Its image is a curve $B\subset X$, and if the degree of $C\to B$ is
$d>0$, the projection formula and Lemma~\ref{surface:degree} give
\[
 -1=K_{X_n}\cdot C=p_n^*K_X\cdot C=d(K_X\cdot B)\geq0,
\]
a contradiction. Thus $X_n$ is minimal. Its class VII numerical
identities and the argument of Lemma~\ref{surface:degree} now give the
asserted canonical basis. No claim about the integral homology of $Y$
enters this argument.
\end{proof}

For reference, the notation used on the cover is summarized here.
\begin{center}
\small
\begin{tabular}{@{}ll@{}}
\hline
$X$, $b=b_2(X)$ & original minimal class VII surface and its second Betti number \\
$Y\xrightarrow{\pi}X$, $T$ & infinite cyclic cover and unit-height deck generator \\
$h$, $K_Y$ & proper height and canonical bundle on $Y$ \\
$X_n=Y/\langle T^n\rangle$ & finite cyclic quotient of degree $n$ over $X$ \\
$P_\lambda$ & flat holomorphic character line bundle on $X$ \\
\hline
\end{tabular}
\end{center}
